% !PRINTED 481
\weberpage{481}

Und ebenso:

6. \gesperrt{Wenn man in einer Invariante von $G$ die Potenzen und Producte der Variablen durch die entsprechenden Ableitungen einer Contravariante ersetzt, so ergiebt sich wieder eine Contravariante.}

\begin{center}
\S. 130.

\gesperrt{Lösung der Gleichung siebenten Grades mit der Gruppe $P_{168}$ durch das Formenproblem der Gruppe $G_{168}$.}
\end{center}

Die lineare Substitutionsgruppe $G_{168}$ hat die bemerkenswerthe Eigenschaft, dass sie mit sich selbst contragredient ist. Denn die erzeugenden Substitutionen $\tau$, $\omega$ von $G_{168}$ (\S.\,115) bleiben durch Transposition ungeändert, und wenn man also irgend eine Substitution der Gruppe transponirt, so erhält man eine Substitution, die gleichfalls in der Gruppe vorkommt.

Die Contravarianten von $G_{168}$ sind also (von der Bezeichnung der Variablen abgesehen) mit ihren Invarianten identisch.

Die in der Gruppe $G_{168}$ enthaltene Octaëdergruppe $G_{24}$, die aus den Substitutionen $\chi^\lambda\omega^\mu\Theta^\nu$ besteht, ist aber von ihrer contragredienten Gruppe verschieden; denn es ist z.\,B. nach \S.\,72, (17)
\[
\Theta^3 = \omega\tau\omega\tau^6,
\]
und die dazu transponirte Substitution
\[
\Theta_1^3 = \tau^6\omega\tau\omega
\]
lässt sich nach \S.\,72, (15) in die Form $\tau\chi^2\Theta^2$ bringen und ist also nicht in $G_{24}$ enthalten.

Es sei nun $\eta_0$ irgend eine zu der Gruppe $P_{24}$ gehörige Function der Grössen $\lambda_0, \lambda_1, \ldots, \lambda_6$ (\S.\,128), z.\,B. die Wurzel $\lambda_0$ selbst. Durch die cyklischen Permutationen $\tau^\nu$ gehe $\eta_0$ in $\eta_0, \eta_1, \ldots, \eta_6$ über. In den Resolventen
\weq{1}{\tilde\omega_r = \sum_{0,6}^{\nu} \varepsilon^{\nu r}\eta_\nu \qquad r = 1, 2, \ldots, 6}
ist dann ein System von Functionen gegeben, die sich durch die Permutationen von $P_{168}$ zwar linear, aber nicht ternär substituiren; da man ja die $\eta_\nu$ selbst linear durch die $\tilde\omega_r$ ausdrücken kann.
